\chapter{Methods}\label{app:methods}
\begin{goals}
\textbf{This appendix} gives every number and drawing wuffle makes as it is computed: the formula, the
conventions, and what happens at the edges. Where wuffle departs from the textbook, it says so. Each
formula here was recomputed independently against wuffle's own code on the Tocher sample and on
synthetic data, and the two agreed.
\end{goals}

The maths lives in \code{@gcu/bearing} and \code{@gcu/stereonet}; wuffle's own modules only arrange
it. Angles are in degrees unless a formula says otherwise.

\section{Conventions}\label{sec:m-conventions}

\subsection{The frame}

A direction is a unit vector $(x, y, z)$ with $x$ east, $y$ north and $z$ \emph{up}. Downward directions
have $z < 0$.

A \textbf{line} with trend $\theta$ and plunge $\varphi$ (plunge positive downward) is
\[ \mathbf{l} = (\cos\varphi \sin\theta,\; \cos\varphi \cos\theta,\; -\sin\varphi). \]

A \textbf{plane} with dip direction $\alpha$ and dip $\delta$ is represented by its \term{pole}, the
downward normal:
\[ \mathbf{p} = (-\sin\delta \sin\alpha,\; -\sin\delta \cos\alpha,\; -\cos\delta), \]
which plunges $90^\circ - \delta$ towards $\alpha + 180^\circ$. The pole of \code{120/40} plunges 50°
towards 300.

Going back, a vector is first folded onto the lower hemisphere (if $z > 0$, every component changes
sign), then
\[ \theta = \mathrm{atan2}(x, y), \qquad \varphi = \arcsin(-z), \]
and for a pole, $\alpha = \mathrm{atan2}(-x, -y)$ and $\delta = \arccos(-z)$, each azimuth taken into
$[0, 360)$.

\subsection{Strike and quadrant notation}

A strike is read by the right-hand rule: the plane dips to the right of it, so the dip direction is
the strike plus 90°, and the rose's strike is the dip direction plus 270°.

An attitude written in quadrant form, such as \code{N30E 40SE}, is a strike and a dip:
\code{N30E} is 030, \code{S40W} is 220, \code{N30W} is 330. The dip's letters choose which of the two
dip directions perpendicular to the strike is meant: the one nearer the letters' azimuth (\code{SE} is
135, so \code{N30E 40SE} is \code{120/40}). Without letters the right-hand rule applies. A direction
written as a number, as in \code{120/40}, is always a dip direction, and letters after its dip are not
used.

\subsection{Axes}

A pole, or a line without a sense of motion, is an \term{axis}: $\mathbf{d}$ and $-\mathbf{d}$ are the
same measurement. Every count, every selection test and the mean treat measurements as axes. A fault's
slip is the one directed quantity (\cref{sec:m-faults}).

\section{Projections}\label{sec:m-projection}

A vector is folded onto the lower hemisphere, then projected. The equal-area (Lambert, Schmidt) and
equal-angle (stereographic, Wulff) projections are
\[ (X, Y)_{\mathrm{area}} = \sqrt{\frac{2}{1 - z}}\,(x, y), \qquad (X, Y)_{\mathrm{angle}} = \frac{1}{1 - z}\,(x, y). \]
For a line plunging $\varphi$ the distance from the centre is
\[ r_{\mathrm{area}} = 2 \sin\frac{90^\circ - \varphi}{2}, \qquad r_{\mathrm{angle}} = \tan\frac{90^\circ - \varphi}{2}, \]
so the primitive circle has radius $\sqrt{2}$ in the first and 1 in the second, and wuffle scales either
to the drawn net. The point sits at azimuth $\theta$ from north, clockwise: $X = r\sin\theta$,
$Y = r\cos\theta$ \parencite{snyder1987}.

\begin{center}
\begin{tikzpicture}[scale=1.6, every node/.style={font=\sffamily\footnotesize}]
  \draw[rule, line width=0.6pt] (0,0) circle (1.414);
  \draw[faint, line width=0.3pt] (-1.55,0) -- (1.55,0) (0,-1.55) -- (0,1.55);
  \node[above, text=muted] at (0,1.55) {N};
  \node[right, text=muted] at (1.55,0) {E};
  % a line 057→30: r = 2 sin 30° = 1, at azimuth 57°
  \draw[amber, line width=0.8pt] (0,0) -- (0.839,0.545);
  \fill[amber] (0.839,0.545) circle (0.045);
  \node[right, text=ink] at (0.88,0.58) {$057\!\to\!30$};
  \draw[ink, ->, line width=0.4pt] (0,0.55) arc (90:33:0.55);
  \node[text=ink] at (0.36,0.70) {$\theta$};
  \node[below right, text=ink] at (0.36,0.25) {$r$};
  \node[text=muted, align=center] at (0,-1.9) {equal-area, lower hemisphere: $r = 2\sin\frac{90^\circ - 30^\circ}{2} = 1$, primitive $\sqrt{2}$};
\end{tikzpicture}
\end{center}

\subsection{The upper hemisphere}

An upper-hemisphere net draws each axis at its upper end, seen from above. wuffle reflects the vector
through the horizontal ($z \to -z$) and projects that, so a line plots at its trend plus 180°, at the
same distance from the centre.

\subsection{Rotation}\label{sec:m-rotation}

A rotated net applies a rotation matrix $\mathbf{R}$ to every direction before projecting it.
\begin{itemize}
\item \code{rotate 120→40} builds the shortest rotation that takes the direction $\mathbf{l}$ of
\code{120→40} to the centre of the net, $(0, 0, -1)$: about the axis $\mathbf{l} \times (0,0,-1)$, by the
angle between them.
\item A drag is an \emph{arcball}: the points under the pointer at the start and end of each movement
are taken back onto the sphere (through the inverse projection; a point outside the primitive is put on
it), and the shortest rotation between them is multiplied onto the current one. The product is
re-orthonormalised after every step, so a long drag never drifts away from a rotation. The matrix is
what the diagram stores.
\end{itemize}
Selections, statistics and the data are unaffected: they live in the unrotated frame, and only the
drawing turns (\cref{ch:rotation}).

\section{The compass}\label{sec:m-compass}

The phone reports its orientation as three angles, $\alpha$ (about the vertical), $\beta$ and $\gamma$,
in the W3C device-orientation convention: the device's frame is reached from east--north--up by turning
$\alpha$ about $z$, then $\beta$ about the new $x$, then $\gamma$ about the new $y$. wuffle builds that
rotation matrix and reads two of its columns:
\begin{itemize}
\item for a \textbf{plane}, the device's $z$ axis (the normal to the screen) is the plane's normal;
\item for a \textbf{line}, the device's $y$ axis (along the long edge, towards the top) is the line.
\end{itemize}
Each is converted as in \cref{sec:m-conventions}, which folds it onto the lower hemisphere. The azimuth
(dip direction or trend) is then corrected for declination, $\mathrm{true} = \mathrm{magnetic} + D$, east
positive.

\subsection{Declination}\label{sec:m-wmm}

$D$ comes from the World Magnetic Model 2025 \parencite{chulliat2025}, evaluated in the page: a
spherical-harmonic expansion of the main field to degree 12, with Schmidt semi-normalised Legendre
functions. The position is converted from geodetic (WGS84, $a = 6378.137$\,km, $f = 1/298.257223563$) to
geocentric coordinates, the field is evaluated there at the reference radius 6371.2\,km, and rotated
back to north, east and down components $X, Y, Z$; then $D = \mathrm{atan2}(Y, X)$. The height is the
GNSS fix's altitude. The coefficients are generated from NOAA's published file, and wuffle's tests pin
the result to the model's twelve published test values. The model is valid from 2025.0 to 2030.0; the
compass doesn't stop using it after that. It recomputes $D$ when the position has moved about a tenth
of a degree (\cref{sec:true-north}).

\section{Contours}\label{sec:m-contours}

Densities are evaluated on a square grid laid over the projected net, 40 nodes a side, each node taken
back to a direction by the inverse projection. Nodes more than about 1\% outside the primitive are left out.
Each contour level is traced by marching squares, and the segments are joined into lines.

\subsection{Kamb}

For $N$ axes and a chosen $k$ (3 unless you change it, \cref{sec:contours}), the counting area is a pair
of opposite caps of half-angle $\vartheta$ with
\[ \cos\vartheta = \frac{N}{N + k^2}. \]
The caps cover the fraction $A = 1 - \cos\vartheta$ of the sphere. For a uniform distribution the count
in them has expectation and standard deviation
\[ E = NA, \qquad \sigma = \sqrt{NA(1 - A)}, \]
and the choice of $\vartheta$ makes $E = k\sigma$ \parencite{kamb1959}. At each node $\mathbf{g}$ wuffle
counts the axes with $|\mathbf{g} \cdot \mathbf{d}| \ge \cos\vartheta$ and stores
\[ \frac{\mathrm{count} - E}{\sigma}, \]
the excess over a uniform distribution, in $\sigma$. The default levels are 2, 4, 6, 8, 10 and 12.

For the Tocher sample ($N = 200$, $k = 3$): $\vartheta = 16.88^\circ$, $A = 0.0431$, $E = 8.61$,
$\sigma = 2.87$.

\begin{note}
Some programs contour the count itself in units of $\sigma$, so that their lowest meaningful level is $k$.
wuffle's level $L$ is their level $k + L$: with $k = 3$, wuffle's 2\,$\sigma$ line is their 5\,$\sigma$.
\end{note}

\subsection{Fisher kernel}

Each axis contributes a Fisher density of concentration $\kappa$ centred on it. The kernel's width is
$s = 90^\circ/\sqrt{N}$ (6.36° for Tocher), and
\[ \kappa = \frac{1}{1 - \cos s}, \qquad f(\mathbf{g}) = \frac{\kappa}{N}\sum_{i=1}^{N} \exp\!\bigl(\kappa\,(|\mathbf{g}\cdot\mathbf{d}_i| - 1)\bigr). \]
For uniformly spread axes the expected value of $f$ is $1 - e^{-\kappa} \approx 1$, so $f$ is in
\emph{multiples of uniform density} (MUD). Averaged over 400 random directions, 200 and 4000 uniform axes
give 0.95 and 0.98. The default levels are 2, 4, 6 and 8.

\subsection{Across the primitive}

Axes near horizontal are counted at both ends, so a cluster on the primitive belongs as much to the
opposite side of the net as to its own. Grid nodes just outside the primitive are therefore taken to
directions just above the horizon, by the same inverse formula continued past the rim; the axial count
there sees the axes near the opposite rim. The contours then run across the primitive and reappear on
the other side, as \textcite{tocher1979} recommends, instead of stopping at the edge.

\subsection{On a rotated or upper-hemisphere net}

On a rotated net the data are rotated and the density evaluated in the rotated frame, so the contours
turn with the points. An upper-hemisphere net is one more transformation of the same kind: the points
are drawn as the lower-hemisphere projection of the z-reflected vector, $(x, y, z) \mapsto (x, y, -z)$,
so the grid is evaluated on the data transformed by the rotation \emph{and} the reflection, one matrix
$Z R$ with $Z = \mathrm{diag}(1, 1, -1)$. The right dihedra are reflected the same way: the field's sign,
$\mathrm{sign}((\mathbf{g}\cdot\mathbf{n})(\mathbf{g}\cdot\mathbf{s}))$, is unchanged when the grid
direction $\mathbf{g}$, the pole $\mathbf{n}$ and the slip $\mathbf{s}$ are all reflected, so reflecting
the faults gives the field at the drawn spot. (Before bearing 0.3.2 the grid had the rotation alone, and
on an upper-hemisphere net the contours sat mirrored through the centre from their points.)

\section{The mean}\label{sec:m-mean}

wuffle's mean is Fisher's \parencite{fisher1953}, made axial. First each axis is turned to the side of the
principal eigenvector $\mathbf{v}_1$ (\cref{sec:m-eigen}): $\mathbf{d}_i \to -\mathbf{d}_i$ when
$\mathbf{d}_i \cdot \mathbf{v}_1 < 0$. Then, with the resultant $\mathbf{S} = \sum \mathbf{d}_i$ and its
length $R$,
\[ \bar{\mathbf{d}} = \mathbf{S}/R, \qquad \bar{R} = R/N, \]
the mean folded onto the lower hemisphere. The concentration is estimated as
\[ \kappa = \frac{N - 2}{N - R} \quad (N \ge 3), \qquad \kappa = \frac{N - 1}{N - R} \quad (N = 2), \]
the first being the small-sample correction of the usual $(N-1)/(N-R)$ \parencite{fisher1987}. When all the
axes coincide ($R = N$), $\kappa$ is infinite. The cone of confidence at level $1 - p$ has half-angle
$\alpha$ with
\[ \cos\alpha = 1 - \frac{N - R}{R}\left[\left(\frac{1}{p}\right)^{1/(N-1)} - 1\right], \]
which at 95\% ($1/p = 20$) is $\alpha_{95}$ \parencite{fisher1953}. It needs $N \ge 2$.

For the Tocher sample: $R = 136.02$, $\kappa = 3.09$, $\alpha_{95} = 6.85^\circ$, and the mean pole gives
the plane 026.7/89.2.

\begin{note}
Turning every axis towards $\mathbf{v}_1$ first is wuffle's way of making Fisher's mean work on axes. On a
cluster it gives the textbook answer. On a girdle, such as Tocher's, the axes have no one side to be
turned to, and the mean and its cone describe little; the eigenvectors are the description to use there
(\cref{sec:mean}).
\end{note}

\section{Eigen analysis}\label{sec:m-eigen}

The orientation tensor of $N$ axes is
\[ \mathbf{T} = \frac{1}{N}\sum_{i=1}^{N} \mathbf{d}_i\,\mathbf{d}_i^{\mathsf{T}}, \]
which is the same for $\mathbf{d}_i$ and $-\mathbf{d}_i$, so the signs of the axes don't matter. Its
eigenvalues $S_1 \ge S_2 \ge S_3$ sum to 1, and its eigenvectors $\mathbf{v}_1, \mathbf{v}_2,
\mathbf{v}_3$ are drawn at their lower ends. From the eigenvalues \parencite{woodcock1977,vollmer1990}:
\[ K = \frac{\ln(S_1/S_2)}{\ln(S_2/S_3)}, \qquad C = \ln\frac{S_1}{S_3}, \]
\[ P = S_1 - S_2, \qquad G = 2(S_2 - S_3), \qquad R = 3S_3, \qquad P + G + R = 1. \]
$K > 1$ is a cluster, $K < 1$ a girdle; $C$ grows with the strength of the fabric. $P$, $G$ and $R$ are
the point, girdle and random parts. Eigen analysis needs three measurements.

For the Tocher sample: $S = 0.548, 0.383, 0.069$; $K = 0.21$, $C = 2.07$; $P = 0.165$, $G = 0.628$,
$R = 0.207$, a girdle.

\subsection{The fabric plots}

The Woodcock diagram places a dataset at $(\ln(S_2/S_3), \ln(S_1/S_2))$, with both axes running from 0 to
7 (larger values are drawn at the edge); the line $K = 1$ is the diagonal. The Vollmer triangle places it
at the weighted mean of its corners, $P$ at the top, $G$ at the lower left and $R$ at the lower right,
with $P$, $G$ and $R$ as the weights (\cref{sec:fabric}).

\begin{center}
\begin{tikzpicture}[scale=3.2, every node/.style={font=\sffamily\footnotesize}]
  \draw[rule, line width=0.6pt] (0,0.866) -- (-0.5,0) -- (0.5,0) -- cycle;
  \node[above, text=muted] at (0,0.866) {P (cluster)};
  \node[below left, text=muted] at (-0.5,0) {G (girdle)};
  \node[below right, text=muted] at (0.5,0) {R (random)};
  % Tocher: P 0.165, G 0.628, R 0.207 → x = −0.5 G + 0.5 R, y = 0.866 P
  \fill[amber] (-0.2105,0.1429) circle (0.022);
  \node[above right, text=ink] at (-0.2105,0.1429) {Tocher};
\end{tikzpicture}
\end{center}

\section{Intersections}\label{sec:m-derived}

The line common to two planes is the cross product of their poles, normalised and folded onto the lower
hemisphere:
\[ \mathbf{b} = \frac{\mathbf{p}_1 \times \mathbf{p}_2}{|\mathbf{p}_1 \times \mathbf{p}_2|}. \]
\begin{itemize}
\item \textbf{Intersection of two datasets} uses each dataset's mean pole (\cref{sec:m-mean}), or the
single plane of a one-row dataset. It reports the angle between the planes, $\arccos|\mathbf{p}_1 \cdot
\mathbf{p}_2|$, and refuses parallel planes.
\item \textbf{The $\beta$ diagram} takes every pair of planes in a dataset: $N(N-1)/2$ lines, each
remembering its pair. Pairs that are parallel to $10^{-9}$ are skipped. A dataset of more than 600 planes
is refused (it would make some 180\,000 lines).
\end{itemize}
For Tocher's 200 planes the $\beta$ diagram has 19\,900 lines.

\section{Selection shapes}\label{sec:m-select}

Every shape is a test on the sphere, applied to a plane's pole or a line, and every test is axial
(\cref{sec:shapes}).
\begin{description}
\item[Cone] of half-angle $\vartheta$ about $\mathbf{a}$: selected when $|\mathbf{d}\cdot\mathbf{a}| \ge
\cos\vartheta$. A plane given as the cone's centre stands for its pole.
\item[Band] of half-width $w$ about the great circle with pole $\mathbf{n}$: selected when
$|\mathbf{d}\cdot\mathbf{n}| \le \sin w$.
\item[Lasso] with vertices $\mathbf{v}_j$ joined by great-circle arcs. The test is made in a
\term{gnomonic projection} centred on $\mathbf{c} = \sum\mathbf{v}_j / |\sum\mathbf{v}_j|$, which
turns great circles into straight lines: with $\mathbf{u}, \mathbf{w}$ two unit vectors perpendicular to
$\mathbf{c}$ and to each other, a direction goes to
\[ \left(\frac{\mathbf{d}\cdot\mathbf{u}}{\mathbf{d}\cdot\mathbf{c}},\; \frac{\mathbf{d}\cdot\mathbf{w}}{\mathbf{d}\cdot\mathbf{c}}\right). \]
An axis is tested at its end on $\mathbf{c}$'s side, by the even--odd rule against the projected polygon.
A lasso wider than a hemisphere selects nothing (none can be drawn on a net).
\end{description}
Because the shapes are on the sphere, the same rows stay selected however the net is rotated or
projected.

\section{Roses}\label{sec:m-rose}

The sectors are centred on multiples of the sector width $w$: an azimuth $a$ falls in sector
$\lfloor (a + w/2)/w \rfloor \bmod (360/w)$, so with 10° sectors, 4.9° counts in the sector centred on
0 and 5.0° in the one centred on 10. A strike is the dip direction plus 270. With \ui{Count both ends}
each azimuth is also counted at $a + 180^\circ$, which a strike rose always does
(\cref{ch:roses}).

The petal of a sector with count $c$, in a rose whose largest count is $c_{\max}$, reaches
\[ r = r_{\max}\sqrt{c/c_{\max}} \quad\mbox{(area-true)}, \qquad r = r_{\max}\,c/c_{\max} \quad\mbox{(linear)}. \]
The area of a sector grows as $r^2$, so the area-true petal's area is proportional to its count.

\section{Faults}\label{sec:m-faults}

\subsection{The sense, and the slip}

A fault row gives a plane (its pole $\mathbf{p}$), a slickenline $\mathbf{l}$ and a sense. Senses are
read as words or single characters:
\begin{center}\small
\begin{tabularx}{\linewidth}{@{}lL@{}}\toprule
\sffamily\bfseries sense & \sffamily\bfseries written as \\\midrule
normal & \code{normal}, \code{n}, \code{2}, \code{-} \\
reverse & \code{reverse}, \code{inverse}, \code{thrust}, \code{i}, \code{r}, \code{1}, \code{+} \\
dextral & \code{dextral}, \code{right-lateral}, \code{d}, \code{3} \\
sinistral & \code{sinistral}, \code{left-lateral}, \code{s}, \code{l}, \code{4} \\
unknown & anything else, \code{unknown}, empty \\
\bottomrule\end{tabularx}
\end{center}
The sense picks the \emph{slip} $\mathbf{s}$, the hanging wall's direction of motion, from $\pm\mathbf{l}$:
\begin{itemize}
\item normal: the one pointing down ($s_z \le 0$); reverse: the one pointing up;
\item dextral: the one opposite the strike vector, sinistral: the one along it. The strike vector is
horizontal, perpendicular to the upward normal $\mathbf{u} = (u_x, u_y, u_z)$, with the plane dipping to
its right: $(-u_y, u_x, 0)/\sqrt{u_x^2 + u_y^2}$. On a fault striking 090 and dipping south, a dextral
slip points west.
\end{itemize}
A vertical plane has two upward normals; wuffle takes the one pointing east (or, for an east--west plane,
north). A fault of unknown sense is drawn, but gets no arrow and is left out of the P and T axes, the
dihedra and the inversion.

\subsection{The slip arrow}

The arrow is drawn at the slickenline's plotted point $\mathbf{q}$ (plunge $\varphi_0$, trend
$\theta_0$) along the horizontal part of the slip, the azimuth $\psi = \mathrm{atan2}(s_x, s_y)$. That
direction splits into a part away from the centre and a part round it,
\[ h_r = \cos(\psi - \theta_0), \qquad h_t = \sin(\psi - \theta_0), \]
and a step of $k$ degrees along the arrow moves the point to plunge $\varphi_0 - h_r k$ and trend
$\theta_0 + h_t k/\cos\varphi_0$. The shaft is 7° long and the head's barbs sit at 55\% of it, 2.4° to
either side. Where the arrow would leave the lower hemisphere ($\varphi_0 - 7 h_r < 0$, pointing outward
near the rim) it is drawn ending at $\mathbf{q}$ instead of starting there. A vertical slip has no
horizontal part and no arrow (\cref{sec:slip-arrows}).

\subsection{P and T axes}

With $\mathbf{u}$ the upward unit normal and $\mathbf{s}$ the unit slip (perpendicular to it),
\[ \mathbf{P} = \frac{\mathbf{u} - \mathbf{s}}{|\mathbf{u} - \mathbf{s}|}, \qquad \mathbf{T} = \frac{\mathbf{u} + \mathbf{s}}{|\mathbf{u} + \mathbf{s}|}, \]
each at 45° to the plane, in the plane that holds the normal and the slip.

\subsection{Right dihedra}\label{sec:m-dihedra}

For each fault $k$, with downward normal $\mathbf{n}_k$ and slip $\mathbf{s}_k$, the function
$2(\mathbf{g}\cdot\mathbf{n}_k)(\mathbf{g}\cdot\mathbf{s}_k)$ is positive in the fault's compressional
quadrants and negative in its extensional ones: $+1$ at its P axis and $-1$ at its T axis. wuffle
averages it over the $m$ faults of known sense,
\[ f(\mathbf{g}) = \frac{1}{m}\sum_{k} 2\,(\mathbf{g}\cdot\mathbf{n}_k)(\mathbf{g}\cdot\mathbf{s}_k), \]
on a 64-node grid, and draws the line $f = 0$ (\cref{sec:dihedra}).

\begin{note}
\textcite{angelier1977} count, at each direction, the faults whose compressional dihedron contains it.
wuffle sums the smooth quadrant function instead, which weighs each fault by how deep in its quadrant
the direction lies. For a single fault the two zero lines are the same, its two nodal planes. For many
faults they are close but not identical, and wuffle's line is the smoother.
\end{note}

The dihedra are computed in the unrotated projection, so they are not drawn on a rotated net.

\subsection{The stress inversion}\label{sec:m-michael}

\textcite{michael1984}'s inversion assumes that each fault slipped along the shear traction the stress
puts on it, and that the shear tractions are of similar size on all the faults. The deviatoric stress
tensor has five unknowns, $\sigma_{11}, \sigma_{12}, \sigma_{13}, \sigma_{22}, \sigma_{23}$, with
$\sigma_{33} = -(\sigma_{11} + \sigma_{22})$. With compression counted positive and the downward normal $\mathbf{n} = (n_1, n_2, n_3)$,
each fault adds the three equations
\[ \sigma\mathbf{n} - (\mathbf{n}\cdot\sigma\mathbf{n})\,\mathbf{n} = \mathbf{s}, \]
linear in the unknowns, to a least-squares system, solved through its normal equations. The principal
axes are the eigenvectors of the result: $\sigma_1$ (the most compressive), $\sigma_2$ and $\sigma_3$,
with the stress ratio
\[ \phi = \frac{\sigma_2 - \sigma_3}{\sigma_1 - \sigma_3}, \qquad 0 \le \phi \le 1. \]
It needs at least four faults of known sense (\cref{sec:stress}).

\begin{why}
On 60 synthetic faults cut through a stress with $\sigma_1$ east--west and horizontal, $\sigma_3$
vertical and $\phi = 0.4$, wuffle's inversion returns $\sigma_1$ at 090→00 and $\sigma_3$ vertical, to a
tenth of a degree, and $\phi = 0.43$. The axes come back exactly; $\phi$ carries the method's assumption
of equal shear stress, which synthetic faults through one tensor don't meet.
\end{why}

\section{The signature}\label{sec:m-signature}

The compass signs the exact bytes of \file{wuffle-log.csv} with ECDSA on the curve P-256 and SHA-256, a
key held in the phone's keystore. The \file{.sig} carries the signature as the 64-byte $r\,\|\,s$ pair and
the public key as SubjectPublicKeyInfo, both in base64 (\cref{sec:format-sig}). wuffle verifies with the
browser's own Web Crypto against the CSV's bytes as they are in the project; the \code{hash} field of the
\file{.sig} is not used.

The key's \term{fingerprint} is the first 8 bytes of SHA-256 over the SubjectPublicKeyInfo, written as 16
hexadecimal digits in groups of four (\code{b2eb 7f10 \ldots}). Two logs with the same fingerprint were
signed by the same key (\cref{sec:field-logs}).
